\chapter{How a MERN Application Works}

Before writing a single line of code, you need a mental model of what actually
happens between a user clicking something in the browser and data changing in
the database. Every feature you build for the rest of this book is just a
variation of the same round trip.

\section{The Full Request Lifecycle}

A MERN stack request always travels through the same set of layers, in the
same order, whether it is a login form or a task creation button.

\begin{diagrambox}[MERN Request Lifecycle]
\begin{verbatim}
 User
  |  clicks / types / submits
  v
 React Component  (JSX + event handler)
  |  calls
  v
 Event Handler  (onClick / onSubmit)
  |  calls
  v
 API function  (services/taskService.js)
  |  calls
  v
 Axios  (src/api/axios.js instance)
  |  sends
  v
 HTTP Request  (method + URL + headers + JSON body)
  |
  v
 Express App  (app.js / server.js)
  |  matches
  v
 Router  (routes/taskRoutes.js)
  |  passes through
  v
 Middleware  (auth check, validation, etc.)
  |  calls
  v
 Controller  (controllers/taskController.js)
  |  calls
  v
 Service / Model layer  (business logic)
  |  calls
  v
 Mongoose  (Task.find / Task.create / ...)
  |  queries
  v
 MongoDB  (actual data on disk)
  |  returns
  v
 Database Response  (documents / write result)
  |  bubbles back up through
  v
 Controller  (shapes { success, data })
  |  sends
  v
 JSON Response  (HTTP status + body)
  |  received by
  v
 Axios  (resolves the promise / rejects on error)
  |  updates
  v
 React State  (useState / context)
  |  triggers
  v
 UI Update  (re-render with new data)
\end{verbatim}
\end{diagrambox}

Each stage has one job:

\begin{itemize}
  \item \textbf{React Component} -- renders UI and listens for user interaction.
  \item \textbf{Event Handler} -- captures the interaction and decides what to call.
  \item \textbf{API function} -- a small wrapper function that knows the exact endpoint and payload shape.
  \item \textbf{Axios} -- the HTTP client; attaches base URL, headers, and credentials.
  \item \textbf{Express App / Router} -- receives the raw HTTP request and matches it to a route.
  \item \textbf{Middleware} -- runs cross-cutting checks (authentication, validation, logging) before the controller.
  \item \textbf{Controller} -- orchestrates the request: validates input, calls the model, builds the response.
  \item \textbf{Mongoose} -- translates JavaScript calls into MongoDB queries and back into JS objects.
  \item \textbf{MongoDB} -- persists and retrieves the actual documents.
  \item \textbf{JSON Response} -- a consistent \texttt{\{ success, data \}} envelope sent back to the client.
  \item \textbf{React State} -- the single source of truth that drives what the UI shows.
\end{itemize}

\begin{notebox}[NOTE]
Every chapter in this book is really just teaching you how to implement one
segment of this diagram correctly and securely. Keep this chain in your head
-- it does not change as the app grows.
\end{notebox}

\section{Concrete Example: ``User Clicks Create Task''}

Let's trace one real interaction end to end in our running example, the
\textbf{Task Management System}.

\begin{diagrambox}[Example: Create Task]
\begin{verbatim}
1. User fills "Title" + "Description", clicks "Create Task"
2. TaskForm.jsx  -> handleSubmit(e) fires
3. taskService.js -> createTask({ title, description })
4. axios instance -> POST /api/tasks  (JSON body, withCredentials)
5. Express router  -> router.post("/tasks", protect, createTask)
6. protect middleware -> verifies JWT cookie, attaches req.user
7. createTask controller -> validates body, calls Task.create({...})
8. Mongoose -> inserts document into "tasks" collection
9. MongoDB -> returns the new document with _id
10. Controller -> res.status(201).json({ success: true, data: task })
11. Axios -> promise resolves with response.data
12. React (setTasks) -> new task appended to state array
13. UI -> TaskList re-renders, new TaskCard appears
\end{verbatim}
\end{diagrambox}

\begin{enumerate}
  \item \textbf{User action:} the user types into two inputs and clicks submit -- no network activity yet.
  \item \textbf{Component:} \texttt{TaskForm.jsx} owns the form state via \texttt{useState} and calls its \texttt{onSubmit} prop.
  \item \textbf{Service function:} \texttt{taskService.js} exports \texttt{createTask()}, the only place that knows the endpoint URL.
  \item \textbf{Axios:} the shared instance adds \texttt{baseURL}, \texttt{Content-Type}, and cookies automatically.
  \item \textbf{Express router:} \texttt{/api/tasks} maps to the \texttt{createTask} controller through the task router.
  \item \textbf{Middleware:} \texttt{protect} rejects the request with 401 if there is no valid session, before any DB access happens.
  \item \textbf{Controller:} validates the payload and delegates persistence to Mongoose.
  \item \textbf{Mongoose / MongoDB:} the model writes a new document and returns it with a generated \texttt{\_id}.
  \item \textbf{Response:} the controller always replies with the same \texttt{\{ success, data \}} shape (more on this in Chapter 8).
  \item \textbf{Frontend update:} the resolved promise's data is pushed into React state, and the list re-renders automatically.
\end{enumerate}

\begin{tipbox}[TIP]
If a bug shows up anywhere in a MERN app, walk this chain step by step --
component, service, axios, router, middleware, controller, model, database --
and check each link. 90\% of bugs are a mismatch between two adjacent links
(wrong URL, wrong field name, missing \texttt{await}, wrong HTTP method).
\end{tipbox}
